%-----------------------------------------------------------------------------%
\chapter*{Kata Pengantar}
%-----------------------------------------------------------------------------%
Puji syukur penulis panjatkan kepada Tuhan Yang Maha Esa atas segala berkat 
dan rahmat-Nya sehingga penulis dapat menyelesaikan skripsi ini dengan 
sebaik-baiknya. 
Selama masa penulisan skripsi ini, penulis telah mendapat banyak dukungan, 
doa, bantuan, inspirasi, dan motivasi dari berbagai pihak. 
Untuk itu, penulis ingin mengucapkan terima kasih kepada:
\begin{enumerate}
    \item Bapak Prof.\ Dr.\ rer.\ nat.\ Hendri Murfi, S.Si., M.Kom.\ selaku 
        Dosen Pembimbing I yang telah membimbing penulis dengan penuh 
        kesabaran dan ketulusan. Arahan dan ilmu yang beliau berikan tidak 
        hanya membantu penulis menyelesaikan penelitian ini, tetapi juga 
        membentuk cara penulis berpikir dalam menghadapi permasalahan 
        secara lebih mendalam.
    \item Ibu Dra.\ Nora Hariadi, M.Si.\ selaku Dosen Pembimbing II yang 
        dengan teliti telah memberikan arahan terkait tata cara penulisan 
        skripsi. Perhatian beliau terhadap setiap detail penulisan sangat 
        membantu penulis dalam menyusun karya ilmiah yang sistematis dan 
        sesuai kaidah akademik.
    \item Bapak Prof.\ Alhadi B., S.Si., M.Kom., Ph.D.\ selaku Dosen Penguji 
        I yang telah memberikan kritik dan saran yang tajam dan membangun, 
        sehingga mendorong penulis untuk berpikir lebih kritis dalam 
        menyempurnakan skripsi ini.
    \item Ibu Dr.\ Nahlia Rakhmawati, M.Si.\ selaku Dosen Penguji II 
        yang telah memberikan masukan dan perspektif yang memperkaya 
        pembahasan dalam skripsi ini.
    \item Ibu Prof.\ Dr.\ Titin Siswantining, D.E.A.\ selaku Dosen 
        Pembimbing Akademis yang senantiasa memberikan arahan, nasihat, 
        dan dukungan selama penulis menjalani masa perkuliahan.
    \item Orang tua penulis, Bapak Timu dan Ibu Yuyun Dharmayanti, yang 
        telah memberikan segalanya. Doa, dukungan, dan kesediaan untuk 
        selalu ada baik di saat semangat maupun di saat lelah menjadi 
        kekuatan terbesar yang menopang penulis hingga sampai di titik ini.
    \item Teman-teman Group DKM (Bagja College) yang telah menjadi bagian 
        dari perjalanan penulis sejak masa SMA dan kebersamaannya terus 
        berlanjut hingga saat ini.
    \item Farhan, Dafa, Nicholas, dan Aditya, teman sepersemesteran yang 
        hadir dari hari-hari pertama kuliah hingga akhir. Kehadiran mereka 
        membuat perjalanan yang panjang ini terasa lebih ringan.
    \item RISTEK Fasilkom UI yang melalui kompetisi \textit{data science}-nya 
        telah membuka jalan bagi penulis menuju dunia yang kini menjadi 
        bagian besar dari kehidupan akademik dan karier penulis.
    \item PT Profina Edvisor yang telah memberikan kepercayaan dan ruang 
        bagi penulis untuk tumbuh sebagai seorang \textit{data scientist}, 
        serta dukungan yang tidak pernah surut selama proses penyelesaian 
        skripsi ini.
    \item Seluruh pihak yang tidak dapat disebutkan satu per satu, yang 
        dengan caranya masing-masing telah turut mengambil peran dalam 
        perjalanan panjang ini.
\end{enumerate}

Penulis menyadari bahwa skripsi ini masih memiliki keterbatasan dan 
kekurangan, baik dari segi penulisan maupun substansi keilmuan. 
Oleh karena itu, kritik dan saran yang membangun sangat penulis harapkan 
sebagai bahan perbaikan ke depan. 
Penulis berharap skripsi ini dapat bermanfaat bagi pembaca dan menjadi 
bagian dari kontribusi nyata dalam pengembangan ilmu pengetahuan, 
khususnya di bidang yang dikaji dalam penelitian ini.

\vspace*{0.5cm}
\noindent Depok, 30 Juni 2026\\[0.1cm]
\penulis